\chapter{Fundamentação técnica: LLMs, agentes e verificação simbólica}
\label{cap:tecnica}

% =====================================================================
% PAPEL DESTE CAPÍTULO (registro técnico):
%   Responde "como isso funciona e por que a combinação é necessária".
%   É a casa canônica de: transformer (2.1.1), inferência probabilística
%   (2.1.2), taxonomia técnica de erros (2.1.3), CoT/PAL/tool-use/
%   self-consistency (2.2), agentes e SMA (2.3), CAS/SymPy (2.4).
%   A hipótese arquitetural é argumentada por inteiro APENAS em 2.5.
%
% NOTA DE REVISÃO: o parágrafo "Descrição: capítulo técnico-central..."
% que abria este capítulo era resíduo do documento de estrutura e foi
% removido. A numeração foi corrigida: "Agentes e sistemas multiagentes"
% passou de 2.2.5/2.2.6 para a Seção 2.3.
% =====================================================================

Este capítulo apresenta os três pilares conceituais sobre os quais se assenta o sistema proposto:
os modelos de linguagem de grande escala, os sistemas multiagentes e a verificação simbólica por
sistemas de computação algébrica. O objetivo não é esgotar cada tema, mas reunir o estritamente
necessário para justificar, ao final, por que a combinação desses elementos é uma necessidade
arquitetural e não uma escolha de conveniência.

\section{Modelos de linguagem de grande escala}
\label{sec:llm}

Os modelos de linguagem de grande escala (\emph{Large Language Models} --- LLMs) são sistemas
computacionais treinados para processar e gerar texto em linguagem natural a partir de padrões
aprendidos em volumes massivos de dados textuais. Compreender seu funcionamento em nível conceitual
é condição necessária para avaliar tanto suas potencialidades quanto suas limitações no contexto
educacional e, em particular, para justificar por que a arquitetura proposta nesta dissertação
requer camadas adicionais de verificação simbólica e avaliação pedagógica.

\subsection{Arquitetura \emph{transformer} em nível conceitual}
\label{sub:transformer}

A quase totalidade dos LLMs contemporâneos é construída sobre a arquitetura \emph{transformer},
introduzida por \citet{vaswani2017attention} no artigo \emph{Attention Is All You Need}. Antes
dessa proposta, os modelos de linguagem predominantes utilizavam redes neurais recorrentes (RNNs),
que processavam sequências de forma estritamente serial --- um elemento de cada vez --- e, por
isso, apresentavam dificuldades em capturar dependências entre palavras muito distantes no texto. A
contribuição central dos \emph{transformers} foi substituir essa lógica sequencial por um mecanismo
de \emph{autoatenção} (\emph{self-attention}), que permite ao modelo avaliar, simultaneamente, a
relevância de cada elemento de uma sequência em relação a todos os demais. Essa capacidade de
processamento paralelo tornou o treinamento de modelos muito maiores viável e deu origem à atual
geração de LLMs.

Para os fins desta dissertação, a arquitetura pode ser compreendida a partir de três conceitos
encadeados. O primeiro é o de \textbf{token}: a unidade mínima de processamento do modelo, que não
corresponde necessariamente a uma palavra inteira, mas a fragmentos de texto determinados
estatisticamente durante o treinamento. A expressão ``função quadrática'', por exemplo, pode ser
segmentada em dois, três ou mais tokens, dependendo do vocabulário do modelo. Crucialmente, o
modelo opera inteiramente sobre sequências de identificadores numéricos associados a esses tokens
--- jamais sobre o texto em si. Isso implica que toda a sofisticação linguística de um LLM é, em
sua raiz, uma transformação de sequências de números.

O segundo conceito é o de \textbf{janela de contexto}: o número máximo de tokens que o modelo pode
considerar em uma única operação de inferência. Tudo que precisa ser levado em conta na geração de
uma resposta --- o enunciado da tarefa, o histórico da conversa, os exemplos fornecidos e a própria
resposta em construção --- deve caber dentro dessa janela. Modelos com janelas mais amplas
conseguem manter coerência em textos mais extensos, o que é relevante para a geração de questões com
contextos elaborados e resoluções detalhadas.

O terceiro conceito é o próprio mecanismo de \textbf{atenção}: o componente pelo qual o modelo
pondera a influência de cada token sobre os demais ao construir sua representação interna do texto.
É esse mecanismo que permite ao modelo reconhecer, por exemplo, que o termo ``raiz'' em um
enunciado matemático pertence ao campo semântico algébrico, e não ao botânico. A atenção opera em
múltiplas camadas e em múltiplas cabeças simultâneas, o que confere ao modelo capacidade de capturar
relações de diferentes naturezas em paralelo.

\subsection{Inferência por previsão probabilística}
\label{sub:inferencia}

Uma vez treinado, o LLM gera texto por meio de um processo denominado \textbf{inferência
autorregressiva}: a cada passo, calcula-se uma distribuição de probabilidade sobre todos os tokens
do vocabulário e seleciona-se o próximo token com base nessa distribuição. O token escolhido é
acrescentado à sequência, que passa a servir de contexto para a previsão seguinte, e o processo se
repete até que um critério de parada seja atingido.

Essa natureza probabilística tem uma consequência fundamental que distingue os LLMs de qualquer
ferramenta de cálculo convencional: \textbf{o modelo não computa, ele prevê}. Quando solicitado a
resolver a equação $2x^2 - 5x + 3 = 0$, o LLM não executa um algoritmo de fatoração; ele gera a
sequência de tokens que, dado seu treinamento, tem maior probabilidade de aparecer após aquele
enunciado. Na maior parte dos casos comuns, esse mecanismo produz respostas corretas, pois o
treinamento expôs o modelo a inúmeros exemplos semelhantes. Porém, ele não oferece nenhuma garantia
de correção formal, o que o diferencia estruturalmente de um sistema de computação algébrica.

Dois parâmetros controlam o grau de variabilidade da geração. A \textbf{temperatura} escala a
distribuição de probabilidade antes da amostragem: valores próximos de zero tornam o modelo mais
determinístico, favorecendo o token mais provável; valores elevados achatam a distribuição,
aumentando a probabilidade de tokens menos prováveis serem selecionados e, por isso, produzindo
saídas mais variadas --- e mais suscetíveis a erros em tarefas que exigem precisão. O parâmetro
\textbf{top-p} (\emph{nucleus sampling}), proposto por \citet{holtzman2020curious}, complementa
esse controle ao restringir o conjunto de candidatos àqueles cuja probabilidade acumulada não
ultrapassa um limiar definido, eliminando da amostragem os tokens de baixíssima probabilidade.

Para aplicações matemáticas --- em que a correção é binária --- a variabilidade introduzida por
esses parâmetros representa um risco concreto. Uma temperatura ligeiramente elevada pode fazer com
que o modelo produza, em diferentes execuções do mesmo \emph{prompt}, respostas distintas e
contraditórias para um mesmo problema. Isso tem implicações diretas sobre a \textbf{reprodutibilidade}
do sistema desenvolvido nesta dissertação, questão retomada no Capítulo~\ref{cap:implementacao}.

\subsection{Limitações conhecidas em raciocínio matemático}
\label{sec:limitacoes}

As limitações dos LLMs em tarefas matemáticas decorrem diretamente de sua arquitetura probabilística
e foram documentadas de forma sistemática na literatura. Três categorias são especialmente
relevantes para o presente trabalho.

A primeira é a \textbf{aritmética e a manipulação simbólica}. Como o modelo opera sobre tokens e não
sobre estruturas algébricas, operações que exigem precisão formal --- fatoração de polinômios,
simplificação de frações algébricas, resolução de sistemas --- não são computadas, mas aproximadas
por padrão estatístico. Quando o padrão é frequente no corpus de treinamento, o resultado tende a
ser correto; quando envolve combinações menos recorrentes de valores ou estruturas, a probabilidade
de erro cresce de forma não previsível.

A segunda limitação é o \textbf{raciocínio encadeado em múltiplos passos}. Em problemas que exigem
uma sequência longa de deduções, cada etapa introduz uma nova fonte de variabilidade probabilística.
Um equívoco em um passo intermediário não apenas invalida aquele passo isolado, mas contamina todos
os subsequentes --- o chamado \emph{cascading error}. A técnica de \emph{chain-of-thought prompting}
(Seção~\ref{sub:cot}) atenua esse problema ao forçar a externalização dos passos intermediários, mas
não o elimina: a cadeia de raciocínio pode ser fluente, coerente na forma e matematicamente
incorreta no conteúdo.

A terceira limitação é a \textbf{sensibilidade ao fraseado do \emph{prompt}}. Pequenas variações na
forma como um problema é enunciado podem produzir respostas estruturalmente diferentes, evidenciando
que o modelo responde ao padrão superficial do texto em vez do conteúdo matemático subjacente. Em um
sistema de geração de questões, essa sensibilidade impõe que os \emph{prompts} dos agentes sejam
rigorosamente formulados e testados.

Em conjunto, essas três limitações fundamentam a hipótese central desta dissertação: nenhum LLM,
utilizado de forma isolada, oferece as garantias necessárias para a geração confiável de questões
matemáticas destinadas ao contexto escolar. A resposta arquitetural a esse diagnóstico --- a
combinação com um verificador simbólico determinístico --- é construída ao longo deste capítulo e
sintetizada na Seção~\ref{sec:hipotese}.

\section{Estratégias de mitigação documentadas na literatura}
\label{sec:mitigacao}

As limitações descritas na seção anterior motivaram uma linha ativa de pesquisa voltada ao
desenvolvimento de técnicas que aumentem a confiabilidade desses modelos sem exigir retreinamento.
As quatro estratégias apresentadas a seguir representam o estado da arte nessa área e constituem o
referencial técnico que orienta as decisões de projeto do sistema desenvolvido nesta dissertação.

\subsection{Chain-of-thought prompting}
\label{sub:cot}

A técnica de \emph{chain-of-thought prompting} (CoT), proposta por \citet{wei2022chain}, consiste em
incluir, nos exemplos fornecidos ao modelo, não apenas o par pergunta-resposta, mas também a cadeia
completa de passos intermediários que conduz à solução. A hipótese subjacente é que, ao ver exemplos
de raciocínio explicitado, o modelo passa a produzir suas próprias cadeias de passos antes de
enunciar a resposta final, o que reduz a probabilidade de saltos inferenciais incorretos.

Os ganhos reportados por \citet{wei2022chain} são expressivos: em \emph{benchmarks} de raciocínio
aritmético e lógico, o CoT superou abordagens de \emph{few-shot} convencional com margem
significativa, especialmente em modelos de maior escala. \citet{wang2023selfconsistency} ampliaram
esse resultado ao demonstrar melhorias adicionais de até 17,9 pontos percentuais no \emph{benchmark}
GSM8K quando o CoT é combinado com amostragem múltipla. Uma variante mais simples, o \emph{zero-shot
CoT}, foi proposta por \citet{kojima2022zero}: basta acrescentar a instrução ``vamos pensar passo a
passo'' ao \emph{prompt} para que o modelo passe a externalizar seu raciocínio, sem necessidade de
exemplos anotados.

Apesar desses avanços, o CoT apresenta limitação relevante para o presente trabalho: ele melhora a
coerência narrativa da resolução, mas não garante correção matemática. O modelo pode produzir uma
cadeia de passos fluente, bem formatada e matematicamente incorreta --- um sinal trocado, uma
fatoração equivocada, um caso desconsiderado ---, e o erro, por estar embebido em texto
aparentemente coerente, tende a passar despercebido sem verificação independente. É precisamente essa
limitação residual que justifica a adição de uma camada de verificação simbólica.

\subsection{Program-aided language models (PAL)}
\label{sub:pal}

Uma abordagem mais radical ao problema da imprecisão matemática é a proposta por \citet{gao2023pal},
denominada \emph{Program-Aided Language Models} (PAL). Em vez de solicitar ao modelo que resolva o
problema diretamente, o PAL instrui o LLM a traduzir o enunciado em linguagem natural para um
programa em Python, cujos passos de execução correspondem às etapas de raciocínio. A solução efetiva
do problema é então \textbf{delegada a um interpretador externo} --- determinístico e exato por
natureza --- que executa o código gerado e retorna o resultado.

Essa divisão de responsabilidades é conceitualmente simples e poderosa: o LLM faz o que faz bem
(compreender linguagem natural e estruturar um plano de solução), e o interpretador faz o que o LLM
não consegue garantir (computar com precisão). Nos experimentos de \citet{gao2023pal}, o PAL superou
modelos muito maiores que utilizavam apenas CoT.

A conexão entre o PAL e a arquitetura desta dissertação é direta e constitui um dos pilares
conceituais do trabalho. O Agente Verificador Simbólico, descrito no Capítulo~\ref{cap:arquitetura}, opera
segundo o mesmo princípio fundamental: em vez de confiar na resolução textual produzida pelo Agente
Gerador, ele delega o cômputo ao SymPy --- um sistema de computação algébrica simbólica --- e compara
o resultado obtido com o gabarito proposto. A diferença em relação ao PAL reside no escopo: onde o
PAL utiliza um interpretador Python de propósito geral, o verificador desta dissertação emprega uma
biblioteca especializada em álgebra simbólica, mais adequada às questões do Ensino Médio e capaz de
lidar com equivalências formais que a aritmética numérica não captura.

\subsection{Tool use e function calling}
\label{sub:tooluse}

Uma generalização da ideia do PAL é o paradigma de \emph{tool use}, que equipa o LLM com acesso a um
conjunto de ferramentas externas --- calculadoras, mecanismos de busca, APIs, interpretadores de
código --- que podem ser invocadas durante a geração da resposta. O modelo aprende não apenas a
produzir texto, mas a decidir \emph{quando} e \emph{como} chamar cada ferramenta e a incorporar os
resultados obtidos ao fluxo de geração. \citet{schick2023toolformer} demonstraram que modelos
treinados com essa capacidade alcançam desempenho competitivo com modelos muito maiores em tarefas
que envolvem aritmética e busca factual. A implementação moderna dessa ideia em APIs comerciais é o
\emph{function calling}: o modelo recebe a descrição de um conjunto de funções disponíveis e pode
emitir chamadas estruturadas a essas funções no meio da geração.

Para o presente trabalho, a distinção relevante é entre \emph{quando usar} e \emph{quando não usar}
essa abordagem. O \emph{tool use} não é a abordagem adotada nesta dissertação: o sistema proposto não
delega ferramentas ao LLM durante a geração, pois isso criaria uma dependência opaca entre o modelo e
a ferramenta, dificultando a rastreabilidade e o controle pedagógico. Em vez disso, a verificação
simbólica é realizada por um agente independente, após a geração, o que preserva a separação de
responsabilidades e permite que o professor visualize tanto a questão gerada quanto o veredicto de
verificação de forma transparente.

\subsection{Self-consistency e amostragem múltipla}
\label{sub:selfconsistency}

A estratégia de \emph{self-consistency}, proposta por \citet{wang2023selfconsistency}, parte de uma
intuição simples: se um problema matemático admite uma resposta correta única, e se o modelo é
executado múltiplas vezes com o mesmo \emph{prompt} (com temperatura maior que zero, para produzir
variações), a resposta mais frequente entre as diferentes execuções tem maior probabilidade de ser a
correta. O processo é análogo a uma votação por maioria entre raciocínios independentes.
\citet{wang2023selfconsistency} reportaram ganhos de 17,9\% no GSM8K e de 11,0\% no SVAMP ao
substituir a decodificação gulosa pela amostragem múltipla com votação por maioria.

A limitação principal é o \textbf{custo computacional}: $k$ execuções do mesmo \emph{prompt} implicam
$k$ vezes o consumo de tokens e, no caso de APIs comerciais, $k$ vezes o custo financeiro. Para
$k=10$, o custo por questão gerada pode tornar-se proibitivo em um sistema destinado a professores da
rede pública com acesso limitado a recursos computacionais. Por esse motivo, o sistema proposto não
adota o \emph{self-consistency} como mecanismo primário de verificação: utiliza-se uma única geração
seguida de verificação simbólica determinística pelo SymPy. Onde o \emph{self-consistency} reduz a
probabilidade de erro por convergência estatística, a verificação simbólica detecta erros de forma
determinística para as classes de questões que o verificador consegue tratar --- mantendo o custo por
questão em patamar compatível com uso educacional real.

Em conjunto, as quatro estratégias formam um mapa das possibilidades disponíveis para aumentar a
confiabilidade de LLMs em Matemática. A Tabela~\ref{tab:estrategias} sintetiza as escolhas feitas
nesta dissertação em relação a cada uma delas.

\begin{table}[htbp]
\centering
\caption{Estratégias de mitigação e seu papel no sistema proposto}
\label{tab:estrategias}
\begin{tabular}{@{}p{3.2cm}p{2.4cm}p{6.5cm}@{}}
\toprule
\textbf{Estratégia} & \textbf{Adoção} & \textbf{Justificativa} \\
\midrule
Chain-of-thought & Parcial & Usado nos \emph{prompts} do Agente Gerador para estruturar a resolução; não suficiente sozinho. \\
PAL / verificação simbólica & Central & Princípio do Agente Verificador Simbólico via SymPy. \\
Tool use & Não adotado & Preservação da separação de responsabilidades e rastreabilidade pedagógica. \\
Self-consistency & Não adotado & Custo computacional incompatível com o contexto de uso; substituído por verificação determinística. \\
\bottomrule
\end{tabular}

\smallskip
{\footnotesize Elaborado pela autora.}
\end{table}

\section{Agentes e sistemas multiagentes}
\label{sec:agentes}

% NOTA DE REVISÃO: esta era a Seção 2.2.5/2.2.6 no rascunho (numerada por
% engano como subseção de "Estratégias de mitigação"). Promovida a 2.3,
% conforme a estrutura prevista.

As estratégias de mitigação apresentadas na seção anterior --- CoT, PAL, \emph{tool use} e
\emph{self-consistency} --- podem ser combinadas em sistemas mais sofisticados, nos quais diferentes
componentes especializados colaboram para resolver tarefas que nenhum deles resolveria adequadamente
de forma isolada. Essa ideia é formalizada pelo paradigma de \textbf{agentes} e \textbf{sistemas
multiagentes}, que constitui a base arquitetural do sistema desenvolvido nesta dissertação.

\subsection{Definição de agente}
\label{sub:agente}

Na formulação clássica de \citet{russell2021aima}, um agente é qualquer entidade que \textbf{percebe}
seu ambiente por meio de sensores e \textbf{age} sobre ele por meio de atuadores. Um agente é
considerado racional quando suas ações são orientadas a maximizar uma medida de desempenho, dado o
conhecimento que possui do ambiente e de suas próprias capacidades. Essa abstração, originalmente
desenvolvida para agentes físicos e sistemas de planejamento clássico, estende-se naturalmente aos
sistemas baseados em LLMs. Um \textbf{agente baseado em LLM} percebe seu ambiente por meio de texto
--- um \emph{prompt} que descreve o contexto, o histórico de ações anteriores e a tarefa corrente ---
e age produzindo uma resposta textual que pode ter efeitos no mundo: chamar uma ferramenta, gerar um
documento, enviar uma mensagem a outro agente. O LLM funciona, nesse modelo, como o módulo de decisão
do agente.

\citet{wang2023agentsurvey} propõem uma taxonomia dos componentes de um agente baseado em LLM
organizada em três dimensões. A primeira é a \textbf{memória}: a capacidade de reter e recuperar
informações ao longo do tempo, de curto prazo (o contexto da conversa atual) ou de longo prazo (um
banco de dados externo consultável). A segunda é o \textbf{acesso a ferramentas}: a possibilidade de
invocar módulos externos que ampliam as capacidades do LLM além da geração de texto. A terceira é a
capacidade de \textbf{raciocínio e planejamento}: decompor tarefas complexas em subproblemas,
executar cada etapa e revisar o plano à luz dos resultados.

Para os fins desta dissertação, cada agente do sistema é definido por uma combinação específica
dessas três dimensões: o Agente Gerador possui acesso ao histórico de tentativas anteriores
(memória); o Agente Verificador possui acesso ao SymPy como ferramenta; e o Agente Crítico Didático
possui uma rubrica pedagógica incorporada ao seu \emph{prompt} como conhecimento estruturado.

\subsection{Propriedades dos sistemas multiagentes}
\label{sub:sma}

Um \textbf{sistema multiagente} (SMA) é composto por dois ou mais agentes que interagem em um
ambiente compartilhado, podendo cooperar, competir ou assumir papéis complementares na resolução de
uma tarefa comum. No contexto dos LLMs, \citet{wu2023autogen} demonstraram que a colaboração entre
múltiplos agentes especializados permite resolver tarefas de maior complexidade do que qualquer
agente individual conseguiria, ao mesmo tempo que melhora a factualidade e a robustez das saídas.
Três propriedades distinguem os SMAs de arquiteturas monolíticas e justificam sua adoção no presente
trabalho.

A primeira é a \textbf{especialização}: cada agente é configurado para executar uma função bem
delimitada, com \emph{prompt} de sistema, ferramentas e critérios de avaliação específicos. Um agente
especializado em verificação simbólica não precisa --- e não deve --- se preocupar com qualidade
pedagógica; um agente especializado em avaliação didática não precisa reproduzir cálculos algébricos.
Essa divisão de responsabilidades reduz a ambiguidade de cada tarefa e melhora a qualidade de cada
componente isoladamente.

A segunda propriedade é a \textbf{verificação independente}: quando um agente produz um resultado que
é avaliado por outro agente com critérios distintos, erros que passariam despercebidos em uma
arquitetura monolítica tendem a ser detectados. Esse princípio é análogo à revisão por pares em
pesquisa científica. No sistema desta dissertação, o Agente Gerador e o Agente Verificador operam com
base em mecanismos completamente diferentes --- geração probabilística \emph{versus} cômputo simbólico
determinístico ---, o que torna a verificação genuinamente independente.

A terceira propriedade é a \textbf{iteratividade controlada}: a presença de um orquestrador que
coordena os ciclos de geração, verificação e revisão permite que o sistema refine progressivamente a
qualidade da questão sem comprometer a rastreabilidade do processo. Cada iteração deixa um registro
auditável, especialmente importante no contexto educacional, em que o professor precisa compreender a
proveniência e o grau de confiança de cada questão produzida.

\subsection{Padrões contemporâneos de orquestração}
\label{sub:orquestracao}

A organização dos agentes em um SMA segue padrões arquiteturais que determinam como as tarefas são
distribuídas e como conflitos são resolvidos. A literatura identifica dois padrões fundamentais e um
ciclo operacional central. O \textbf{padrão supervisor} organiza os agentes em uma hierarquia com um
coordenador que decide quais agentes acionar, em que ordem e com quais entradas, oferecendo controle
centralizado e rastreabilidade clara. O \textbf{padrão roteador} direciona cada tarefa ao agente mais
adequado com base em características da entrada, sendo mais apropriado quando as tarefas são
heterogêneas, mas oferecendo menor controle sobre ciclos de revisão. Os \textbf{ciclos de revisão}
(geração--crítica--revisão) combinam os dois anteriores: um agente gera uma proposta, um segundo a
avalia criticamente, e o resultado é devolvido ao gerador para refinamento, em um ciclo que se repete
até que um critério de parada seja satisfeito. Esse último padrão é o mecanismo central do sistema
proposto nesta dissertação.

Três \emph{frameworks} de código aberto dominam atualmente a implementação de SMAs baseados em LLMs.
O \textbf{AutoGen} \citep{wu2023autogen}, da Microsoft Research, organiza os agentes como entidades
conversáveis que trocam mensagens, o que facilita a implementação de ciclos de revisão, mas pode
tornar o fluxo de controle menos previsível. O \textbf{LangGraph} baseia-se em uma representação
explícita do fluxo de execução como um grafo dirigido, oferecendo controle preciso e facilitando a
depuração --- adequado a sistemas com lógica de decisão complexa. O \textbf{CrewAI} adota uma metáfora
de equipe de trabalho, com papéis definidos declarativamente, o que facilita a prototipação rápida,
embora ofereça menos controle granular. A Tabela~\ref{tab:frameworks} sintetiza essas características.

\begin{table}[htbp]
\centering
\caption{Comparativo de \emph{frameworks} para sistemas multiagentes}
\label{tab:frameworks}
\begin{tabular}{@{}p{2.2cm}p{3.0cm}p{2.8cm}p{4.0cm}@{}}
\toprule
\textbf{Framework} & \textbf{Abstração central} & \textbf{Controle de fluxo} & \textbf{Adequação ao sistema proposto} \\
\midrule
AutoGen & Conversa entre agentes & Implícito (mensagens) & Adequado para prototipação; controle limitado. \\
LangGraph & Grafo de execução & Explícito (grafo) & Alta aderência; controle preciso sobre ciclos. \\
CrewAI & Equipe com papéis & Declarativo (alto nível) & Rápida configuração; menor controle granular. \\
\bottomrule
\end{tabular}

\smallskip
{\footnotesize Elaborado pela autora.}
\end{table}

A decisão sobre adotar ou não um desses \emph{frameworks} --- e qual --- é uma escolha de projeto,
apresentada e justificada na Seção~\ref{sec:decisoes}, à luz dos requisitos do sistema.

\section{Sistemas de computação algébrica como verificadores}
\label{sec:cas}

A discussão das seções anteriores convergiu para um diagnóstico claro: LLMs geram texto
matematicamente plausível, mas não executam cômputos com garantia de correção formal. A estratégia de
verificação simbólica proposta nesta dissertação responde a esse diagnóstico com uma classe de
ferramentas que existe precisamente para isso: os \textbf{sistemas de computação algébrica}
(\emph{Computer Algebra Systems} --- CAS).

\subsection{Computação algébrica simbólica}
\label{sub:cas-conceito}

Um sistema de computação algébrica é um programa capaz de manipular expressões matemáticas em sua
forma simbólica --- isto é, de operar sobre símbolos e estruturas algébricas exatas, sem recorrer a
aproximações numéricas. A distinção em relação ao cálculo numérico é fundamental: enquanto uma
biblioteca numérica representa $\sqrt{2}$ por sua aproximação decimal ($1{,}41421356\ldots$), um CAS
mantém $\sqrt{2}$ como objeto simbólico exato, capaz de participar de operações como
$\sqrt{2}\cdot\sqrt{2}=2$ sem qualquer perda de precisão. Essa distinção tem consequências práticas
diretas: confirmar que as raízes de $x^2-2=0$ são $x=\pm\sqrt{2}$, ou que $\frac{x^2-1}{x-1}=x+1$ para
$x\neq 1$, exige manipulação simbólica, não avaliação numérica em pontos isolados.

Do ponto de vista histórico, a computação algébrica simbólica tem suas raízes no projeto MACSYMA,
desenvolvido no MIT na década de 1960. Nas décadas seguintes, o Maple e o Mathematica tornaram-se
referência, mas são proprietários e de licenciamento custoso, o que limita sua adoção em contextos
educacionais de acesso restrito. Esse cenário motivou o desenvolvimento do \textbf{SymPy}: uma
biblioteca de computação algébrica simbólica escrita inteiramente em Python, de código aberto e
licença permissiva (BSD), que não depende de nenhum sistema externo e pode ser instalada e utilizada
gratuitamente em qualquer ambiente computacional \citep{meurer2017sympy}. Por essas características
--- abertura, integração nativa com Python e custo zero ---, o SymPy foi adotado como verificador
simbólico no sistema proposto.

\subsection{SymPy: capacidades relevantes para o Ensino Médio}
\label{sub:sympy}

Interessa detalhar as capacidades do SymPy diretamente relacionadas aos conteúdos matemáticos do
Ensino Médio que o sistema precisará verificar.

\textbf{Resolução de equações e sistemas.} A função \texttt{solve} recebe uma expressão ou lista de
expressões e retorna o conjunto-solução exato, incluindo radicais e números complexos quando
pertinente. Para sistemas, retorna as soluções como dicionários de atribuições, permitindo comparação
direta com o gabarito proposto pelo Agente Gerador.

\textbf{Manipulação de polinômios e expressões algébricas.} Operações como expansão
(\texttt{expand}), fatoração (\texttt{factor}), simplificação (\texttt{simplify}) e redução ao mesmo
denominador (\texttt{together}) são executadas de forma exata, permitindo verificar se a forma
fatorada de um gabarito é equivalente ao polinômio original.

\textbf{Cálculo de derivadas, integrais e limites.} Para conteúdos de introdução ao Cálculo, o SymPy
oferece \texttt{diff}, \texttt{integrate} e \texttt{limit}, que operam simbolicamente e retornam
resultados exatos.

\textbf{Trigonometria simbólica.} As funções trigonométricas simbólicas, com \texttt{trigsimp} e
\texttt{expand\_trig}, permitem verificar a equivalência entre expressões distintas que representam o
mesmo valor, como $\sin^2\theta + \cos^2\theta = 1$ ou $\cos(2\theta) = 1 - 2\sin^2\theta$.

\textbf{Equivalência simbólica.} A capacidade mais crítica para a verificação de gabaritos é
determinar se duas expressões $A$ e $B$ são iguais, o que se realiza aplicando \texttt{simplify} à
diferença $A-B$: se o resultado for zero, as expressões são equivalentes. Essa abordagem captura
identidades algébricas gerais, não apenas coincidências em valores particulares.

\subsection{Limites da verificação simbólica}
\label{sub:cas-limites}

A despeito de suas capacidades, a verificação simbólica via SymPy apresenta limites que o sistema
precisa reconhecer e tratar explicitamente. O primeiro diz respeito a \textbf{problemas de geometria
com construção}: questões que envolvem demonstração de propriedades, construção com régua e compasso
ou argumentação sobre figuras não se prestam à verificação simbólica direta --- não há conjunto-solução
a comparar, mas uma argumentação lógica. Para esse tipo de questão, o Agente Verificador emite o
veredicto \emph{não-verificável} e devolve a validação ao professor.

O segundo limite refere-se a \textbf{enunciados em linguagem natural ambíguos}: o SymPy verifica
expressões formalizadas, não prosa. Quando o Agente Gerador produz um enunciado ambíguo, o
Verificador pode formalizar uma interpretação diferente da pretendida, produzindo um veredicto
incorreto --- risco mitigado pelo Agente Crítico Didático, que avalia a clareza do enunciado.

O terceiro limite é a \textbf{incompletude dos algoritmos de simplificação}: para certas classes de
expressões, \texttt{simplify} pode não reduzir $A-B$ a zero mesmo quando as expressões são
equivalentes. Para tratar esse caso, o sistema adota uma \textbf{verificação numérica complementar}:
avalia as expressões em um conjunto de pontos aleatórios e compara os resultados com tolerância
definida. Essa verificação não oferece prova de equivalência, mas constitui evidência probabilística
forte --- se duas expressões concordam em $k$ pontos aleatórios independentes, a probabilidade de
discordância em algum ponto não testado decresce exponencialmente com $k$. O veredicto resultante é
registrado como \emph{aprovado com ressalva numérica}.

\section{Síntese: a hipótese arquitetural desta dissertação}
\label{sec:hipotese}

% CASA CANÔNICA da hipótese arquitetural: este é o ÚNICO lugar onde o
% argumento completo é desenvolvido. Os capítulos anteriores apenas o
% anunciam (gancho) e referenciam esta seção.

Os pilares conceituais apresentados neste capítulo não são contribuições isoladas. Eles se articulam
em uma hipótese de trabalho que pode ser enunciada de forma direta: \textbf{a combinação de um LLM
gerador, um verificador simbólico determinístico, um agente avaliador pedagógico e um orquestrador de
ciclos de revisão produz questões matemáticas de maior confiabilidade do que qualquer um desses
componentes poderia produzir de forma isolada.}

A lógica dessa combinação parte de uma constatação sobre limites complementares. O LLM é o único
componente capaz de compreender uma especificação pedagógica em linguagem natural --- tópico,
habilidade BNCC, nível cognitivo, contexto temático --- e produzir um enunciado plausível, uma
resolução passo a passo e, no caso de múltipla escolha, distratores informativos. Nenhum sistema
determinístico realiza essa tarefa. Mas essa mesma flexibilidade probabilística é a fonte do
\emph{erro silencioso} identificado na Seção~\ref{sec:erro-silencioso}: o LLM erra com fluência, sem
sinalizar incerteza.

É nesse ponto que o SymPy intervém. O verificador simbólico não é um segundo LLM revisando o trabalho
do primeiro: ele opera segundo uma lógica completamente distinta, baseada em transformações algébricas
determinísticas. Por isso, seus erros --- quando ocorrem --- são de natureza estruturalmente diferente
dos erros do LLM, o que torna a probabilidade de falha simultânea dos dois componentes na mesma
direção muito menor do que em qualquer arquitetura homogênea. Como já discutido a propósito do PAL
(Seção~\ref{sub:pal}) e do \emph{self-consistency} (Seção~\ref{sub:selfconsistency}), a verificação
independente por mecanismo diferente é mais do que uma redundância: é uma garantia qualitativa que a
simples repetição do LLM não ofereceria.

A correção matemática, porém, não esgota os critérios de qualidade de uma questão. Uma questão pode
estar algebricamente impecável e ser pedagogicamente inadequada --- enunciado ambíguo, nível cognitivo
incompatível com a habilidade declarada, distratores sem função formativa. Essa dimensão é coberta
pelo \textbf{agente crítico didático}, que avalia a questão com base em uma rubrica derivada dos
referenciais da literatura de educação matemática. É precisamente por ser externo que esse critério
tem valor: a rubrica é construída independentemente da questão julgada.

O \textbf{orquestrador} articula esses três componentes em ciclos de revisão controlados. Ele não
gera, não verifica e não avalia: decide. Com base nos veredictos do verificador e do crítico,
determina se a questão está aprovada, se deve retornar ao gerador com instruções de melhoria ou se
deve ser descartada por inviabilidade de correção dentro do número máximo de iterações. Essa política
de decisão é o que transforma a arquitetura de uma \emph{pipeline} linear em um sistema adaptativo.

O resultado dessa combinação não é a eliminação do erro --- nenhum sistema automatizado pode garantir
isso em linguagem natural --- mas a sua \textbf{visibilidade e rastreabilidade}. O professor recebe não
apenas a questão, mas os veredictos estruturados de cada agente, o que posiciona o sistema como
ferramenta de apoio à autoria docente, e não como substituto do julgamento profissional. Para que essa
arquitetura funcione, contudo, a rubrica do agente crítico precisa estar ancorada em referenciais
pedagógicos reconhecidos. O Capítulo~\ref{cap:didatica} cumpre essa função, apresentando os referenciais
de elaboração de questões matemáticas e consolidando-os em critérios de qualidade operacionalizáveis.
